\chapter*{Glosario}
\addcontentsline{toc}{chapter}{\hspace{1em} Glosario}

\textbf{BEACON:} mensaje periódico con el que un nodo anuncia su identidad y su posición a quienes lo oyen directamente; no se reenvía.\\

\textbf{COBERTURA MEDIA:} promedio en el tiempo de la fracción de Nodos de usuario vivos que comparten componente de radio con algún Gateway vivo.\\

\textbf{ESTIRAMIENTO:} saltos que recorrió un paquete entregado divididos por los del camino más corto por radio en el instante del envío.\\

\textbf{GATEWAY:} nodo de la malla que retransmite tráfico y extiende la cobertura; es el rol que puede moverse para reubicarse.\\

\textbf{INUNDACIÓN:} protocolo de comparación sin tablas de rutas en el que cada nodo difunde una vez cada paquete nuevo.\\

\textbf{LOTE:} archivo que define un conjunto de corridas (escenarios por semillas) y del que se obtiene un resumen estadístico.\\

\textbf{NODO CRÍTICO:} punto de articulación de la malla: si cae, la red queda dividida en más componentes.\\

\textbf{NODO DE USUARIO:} nodo final que necesita cobertura de la malla.\\

\textbf{PLANO DE DATOS:} parte del simulador que lleva el tráfico de usuario salto a salto usando la tabla de rutas de cada nodo.\\

\textbf{RED DE EXPANSIÓN DE COBERTURA:} malla de Gateway desplegada para llevar conectividad a Nodos de usuario en un área, sin coordinador central.\\

\textbf{RUTA FANTASMA:} ruta vigente en la tabla de un nodo hacia un destino con el que ya no existe ningún camino por radio.\\

\textbf{RUTA VIGENTE:} ruta presente en la tabla y refrescada por un OGM hace menos que el tiempo límite.\\

\textbf{SEMILLA:} valor inicial del generador de números aleatorios; con la misma semilla y configuración se repite la misma corrida.
